\documentclass[UTF8,11pt]{ctexart}
\usepackage{amsmath,amssymb,amsthm,mathtools}
\usepackage{booktabs,longtable,array,geometry,enumitem}
\usepackage{xcolor}
\definecolor{linkblue}{rgb}{0.10,0.20,0.45}
\usepackage[colorlinks=true,linkcolor=linkblue,citecolor=linkblue,urlcolor=linkblue]{hyperref}
\geometry{a4paper,margin=2.4cm}
\setlength{\parskip}{4pt}
\setlength{\emergencystretch}{3em}
\sloppy
\newcommand{\Lam}{\Lambda}
\newcommand{\R}{\mathbb{R}}
\newcommand{\E}{\mathbb{E}}
\newcommand{\M}{\mathcal{M}}
\newcommand{\dd}{\,\mathrm{d}}
\newcommand{\tr}{\operatorname{tr}}
\DeclareMathOperator{\diag}{diag}
\DeclareMathOperator{\supp}{supp}
\newcommand{\cstar}{c_{\mathrm{w}}}
\newcommand{\file}[1]{\texttt{#1}}
\newcommand{\ok}{\textcolor[rgb]{0.05,0.45,0.15}{\bfseries [已证]}}
\newcommand{\num}{\textcolor[rgb]{0.55,0.35,0.00}{\bfseries [数值]}}
\newcommand{\todo}{\textcolor[rgb]{0.70,0.15,0.10}{\bfseries [未完成]}}
\title{\bfseries 前三大特征值之和的上界\\[4pt]
\large 论文优化说明、方法与实质性进展总结、排版建议及汇报讲稿}
\author{}
\date{2026 年 11 月 2 日}
\begin{document}
\maketitle
\vspace{-1.5em}

\begin{abstract}
\noindent 本文件配合 26 页自编译汇报幻灯片（\file{汇报PPT\_自编译\_前三大特征值上界\_20261102.pptx/.pdf}）使用，内容分四块。
第一块记录本轮按优先级对论文正文 \file{top3\_joint\_moment\_bound.tex} 所做的加固，
即审核报告中两处承重缺口 G1、G2 的修补与七条表述建议 S1--S7 的落实，并新增附录 A 交代全部常数的来源与机器验证协议。
第二块把论文所用的方法完整复述一遍：结构线（二阶变分 $\to$ 符号类秩控制 $\to$ 矩保持变分 $\to$ 分离二元观测 $\to$ 极小依赖）
与证书线（联合三、四阶矩恒等式 $\to$ 带符号四次控制多项式 $\to$ 八项非负分解），以及它为什么能超过 STW 的 $(3+\sqrt5)/4$。
第三块给出实质性进展的量化：改进量 $0.0243437$、冠军律处证书松驰的七项归属、本框架的天花板
$\frac{29+10\sqrt5}{40}=1.2840170$、以及数值上确界 $\rho_0=\frac{9+\sqrt5+2\sqrt{35}}{18}=1.2815682$；
同时如实列出当前阻碍。第四块是逐页讲稿与论文排版建议清单，可直接用于组会汇报。
\end{abstract}

\vspace{-0.5em}
\noindent\textbf{一句话结论。}
主定理 $\Lam_3(G)\le cn-3$（$c=1.284673255723918505\ldots$）\ok，其证明中两处跳步已补齐、经独立审核与机器复验；
下一步的主攻方向是证明``小分支为空''，把 $c$ 推到框架天花板 $\frac{29+10\sqrt5}{40}$，
精确上界 $\M=\rho_0$ 作为长期目标，不宜在答辩前承诺。

\tableofcontents
\newpage

\section{本轮完成的工作}

\subsection{论文加固（优先级 P0，已完成）}

审核报告（\file{审核报告\_top3\_joint\_moment\_bound.pdf}，33 页）给出的结论是：主定理成立，
但证明链上有两处\emph{承重}的跳步，必须补上才能投稿或答辩。本轮已把这两处补进正文，
并顺带落实了七条表述建议。所有改动都以最小侵入为原则，不改动任何数学内容的实质，
论文重新用 \file{pdflatex} 编译两遍，从 19 页变为 20 页；其后第二次查新复核（补引与注 3.2）扩为 22 页纠正版，均无警告、无未解析引用。

\textbf{G1（引理 2.6，\file{contact-structure}，阈值公式）。}
原文在论证``极大行在零集意义下唯一''时，只写了``Uniqueness forces $\int_A F=0$''一句，
没有说明这个唯一性从何而来：如果 $A_0=\{v\cdot F=0\}$ 上存在两个不同的极大行，
它们的差在 $A_0$ 上 supported，而阈值公式恰恰要对\emph{每一个}原始行成立，
所以这一步不是修辞，而是承重的。修补办法是把 $A_0$ 上的候选行与已知行 $v$ 做\emph{配对}：
考虑 $h_v+\mathbf 1_A$（$A\subset A_0$ 可测），它与 $h_v$ 有相同的支撑取值，
于是二者都必须是同一变分问题的极大元；由严格凸性（或此处可直接由椭球恒等式）
得 $\int_A F=0$ 对一切这样的 $A$ 成立，故 $F=0$ a.e.\ on $A_0$，
再由椭球恒等式知 $A_0$ 零测。正文中这一处现在写成约十行的完整论证，
并把``$A_0$ 零测''与``极大行唯一''拆成两个独立结论，
避免审稿人在同一个位置提出同样的质疑。

\textbf{G2（引理 4.2，\file{binary}，分离观测引理）。}
$0\in H$、$\dim H=2$ 的那一支，原文只论证了``对任何在 $I$ 上非常值的观测 $b$，
$I_1$ 与 $I_0$ 都至少含三点，从而 $s\ge6$，矛盾''，
却漏掉了``所有观测在 $I$ 上都是常值''这一情形本身需要被排除或单独处理。
修补只需一句话：若每个分离观测在 $I$ 上常值，则 $I$ 上的原子在该观测族下不可区分，
与 $I$ 取自正团且 $|I|\ge2$ 的选取矛盾，故该情形不出现。
这一句撑起 $m\le24$，而 $m\le24$ 又撑起最终证书中 $\chi-\frac{C^2}{24}\ge0$ 一项，
所以它虽然只有一句话，位置却很关键。

\textbf{S1--S7（表述与可读性）。}
S1 在引理 2.3 处补一句：块均值零子空间使 $\Sigma_3$ 等于商矩阵的谱和，
这样 $\rho_0$ 的下界构造才与``六块阶梯核，对角块取 1''的写法自洽。
S2 在引理 2.5 处显式写出 $A\sqcup B$ 试探标架的单位化归一，
原文``direct-sum identification''一带而过。S3 在 \S5.1 补半句``$R_t=Y$ 故 $\E[R_t]=x$''，
使 $L_q-L=q_0(\zeta-x)^2(\zeta+2x)/x^2$ 的推导不留空档。
S4 把引理 8.2 中常数 $3/2$ 的非紧性写成一条注（\file{rem-theta}）：
真实上确界 $z_*^2/\sigma$ 在两个分支上分别是 $1.200224$ 与 $0.615384$，
所以 $3/2$ 有明确余量、不是随手取的。S5 新增附录 A（\file{sec-appendix}），
把论文里出现的整数与有理常数分成``由插值条件唯一确定''与``先选定再用精确算术验证''两类，
逐条列明，并写下验证协议：所有恒等式、配方与``$>0$''断言都在精确有理算术与二次域
$\mathbb Q(\sqrt3)$、$\mathbb Q(\sqrt5)$ 中复核，从不用浮点；
点态不等式额外在 $2\times10^6+1$ 个网格点上扫描；有限图转移在 $3\le n\le6$ 的全部
$33864$ 个标号图上检查。S6 把定理 1.2 重排成 (a)--(d) 四条，
``至多 24 个支撑点''与``原 graphon 几乎处处等于二值阶梯核''不再混在一句话里。
S7 新增一条注（\file{rem-radii}）说明引理 6.3 的 $r_+/r_-\le3$ 与 \S9 的 $32/27$ 阈值相容，
并指出把 $4/3$ 换成 $32/27$ 可带来 $12.5\%$ 的额外收益，但当前证明不需要，故不改。

\subsection{汇报材料（本轮新增，自编译）}

幻灯片 26 页，不用任何模板 skill：版式由 \file{python-pptx} 直接排（16:9、统一色板、统一页脚与页码），
公式（STW 界、$c$ 的闭式、律界、天花板、$\rho_0$、以及 $D$ 与 $d$）与三张图表
（松驰七组归属、常数阶梯瀑布、八项对数条形）由 \file{matplotlib} 渲染为 PNG 后嵌入，
PDF 版经 WPS 的 COM 接口（\file{KWPP.Application} 的 \file{SaveAs(...,32)}）导出。
配色与标签沿用审核报告的约定：绿＝已证、蓝＝已核或方法、橙＝数值或非紧常数、红＝未完成或阻碍。
本文件即配套讲稿，\S6 的逐页口播与 26 页一一对应。
上一版基于 dashi 模板的 15 页幻灯片已弃用，留档于 \file{\_dashi\_弃用/}。
全部产物与论文 \file{.tex}/\file{.pdf} 一并推送到 git 仓库
\file{github.com:jhy28740-jhy/graph-theory-study.git} 的 \file{main} 分支。

\section{从图极限到本文：论文是如何写出来的}

本节面向只读过 Nikiforov 图极限原文的读者，交代全文的写作逻辑：先搭一座桥，再分两条线推进。

\subsection{底座：graphon 与紧性}
对称可测 $W:[0,1]^2\to[0,1]$；有限图 $G$ 对应阶梯 graphon $W_G$；blow-up 不改变 graphon；
割度量给出紧性，使``在某类图中取极值''可以 pass 到极限 graphon 上取到极大值。
Nikiforov 的解析方法即：把谱极值量写成 graphon 泛函，取极限得极值 graphon，再作变分刻画。
本文的第一个难点是 $\lambda_2,\lambda_3$ 不是同态密度，不能直接套密度型紧性。

\subsection{第一块砖：Ky Fan 桥}
Ky Fan 最大原理 $\lambda_1+\lambda_2+\lambda_3=\max_{\dim U=3}\operatorname{tr}(P_U A)$
把 $\Lambda_3$ 变成迹型泛函，在 graphon 极限下可转移。极值 graphon 的三维最大化子空间给出
特征函数 $f_1,f_2,f_3$；点 $x\mapsto v(x)=(f_1,f_2,f_3)(x)\in\mathbb R^3$；
结构被律 $\mu=\mathrm{law}(v(x))$ 捕获，归一化后各向同性 $\mathbb E[vv^\top]\propto I$。
于是谱问题约化为：对所有各向同性律控制一个由 $\mu$ 决定的量（毕业论文 Lemma 2.3）。

\subsection{结构线（\S2--\S4）：极值律的形状}
\S2 变分结构把 $\mu$ 约束进一个刚性家族；\S3 用矩关系的传播与符号区域（半平面交、直线排列）
证明支撑有限；\S4 用二值观测分离得支撑 $\le 24$ 原子、正点团 $\le 6$、且二值。
结构头条：每个极大 graphon 的典范律至多 24 个原子。

\subsection{证书线（\S5--\S9）：把上界变成可机检不等式}
\S5 矩与回归把谱不等式翻译成少数几个矩的不等式；\S6 的带号四次型与联合缺陷恒等式使其
非正性等价于目标上界；\S7 方向与径向分开控制；\S8 用无三次项的接触多项式证明四次松驰中的
储备非负且有余，组装完整联合证书；\S9 汇总得 $c=1.2846732557\ldots$，并转移回有限图、
给出精确有限修正 $-3$，即 $\Lambda_3\le cn-3$。

\subsection{为什么赢 STW、以及还差什么}
STW 的 $(3+\sqrt5)/4\approx1.30902$ 来自较粗的二次放缩；本文用联合矩$+$四次松驰$+$储备，
把 24-原子结构里更多信息榨出来，得 $c=1.284673$，改进 $0.024344$（约 1.86\%）。
但 $c$ 与加权/吹起常数 $c_w=1.281568$ 之间仍差 $0.003105$，最优常数与极图仍开放。

\section{论文所用的方法}

这一节把论文的技术路线完整复述一遍，目的是让导师在不读正文的情况下，
能够判断``改进从哪里来''以及``哪里还是软的''。为方便对照，
凡引用正文结论都给出其 \file{label} 名与审核报告中的编号。

\subsection{问题设定与已知基准}

设 $G$ 是 $n\ge3$ 阶有限简单无向图，$\lambda_1(G)\ge\lambda_2(G)\ge\cdots$ 为邻接矩阵特征值，
按代数序排列，\emph{保留符号}，记 $\Lam_3(G)=\lambda_1+\lambda_2+\lambda_3$。
问题是：只给定 $n$，$\Lam_3(G)$ 最大能到多少。
这个问题在 Ebrahimi 的硕士论文中已经出现，Mohar 给出了一般的前 $k$ 大特征值之和的界；
最近的进展是 Sivashankar--Tang--Wakhare（下称 STW）命题 6.1：
\begin{equation}\label{eq:stw}
 \Lam_3(G)\le\frac{3+\sqrt5}{4}\,n-3 .
\end{equation}
他们的证明走的是秩三绝对投影常数 $\gamma(3)=\frac{1+\sqrt5}{2}=\varphi$
（所需的投影常数界由 Der\k{e}gowska--Lewandowska 给出）。

关键的归约是：若 $X\in\R^3$ 满足 $\E XX^T=I_3$，$X'$ 是独立副本，令
\[
 A=\E|X\cdot X'|,\qquad B=|\E X|^2 ,
\]
则\emph{分别}估计 $A\le\varphi$ 与 $B\le1$ 就得到 \eqref{eq:stw}。
也就是说，STW 的界来自两次独立放缩的和。要想改进，
唯一可能的来源是对 $A$ 与 $B$ \emph{同时}取大值的定量限制——
这个``混合投影常数''的问题 STW 在其 \S6.2 已经提出但没有解决。
论文直接采用它的各向同性律表述，这是全篇的出发点，
也解释了为什么本文的改进是\emph{结构性}的而不是把某个常数调紧一点：
在 $A$、$B$ 分别取等的方向上，$\gamma(3)=\varphi$ 已经是精确值，没有任何余量可捞。

\subsection{主定理与常数的闭式}

定义
\begin{equation}\label{eq:c}
 c=\frac{106106902139519}{107310156800000}+\frac{41\sqrt3}{240}-\frac{53}{9000000}
   =1.284673255723918505\ldots .
\end{equation}
\textbf{定理 1.1（\file{main}）}\ok\ 对一切 $n\ge3$ 阶有限简单图，$\Lam_3(G)\le cn-3$；
并且对一切满足 $\E XX^T=I_3$ 的 $\R^3$ 上 Borel 概率律，
\begin{equation}\label{eq:lawbound}
 \E|X\cdot X'|+|\E X|^2\le 2c ,
\end{equation}
只假设二阶矩有限，允许零向量与无界半径。
系数较 \eqref{eq:stw} 降低 $0.024343738651028919\ldots$，常数项 $-3$ 不变。

\eqref{eq:lawbound} 是真正的主体，图论版本只是它的转移：
论文用 graphon 语言定义
$\Sigma_k(W)=\sup_{\E FF^T=I_k}\iint W(x,y)F(x)\cdot F(y)\dd x\dd y$、
$\M=\sup_W\Sigma_3(W)$，并说明 $\M$ 可达；
再由``有限图的邻接矩阵归一化后是一个 step graphon''把 \eqref{eq:lawbound} 转成图界，
转移过程中产生精确的 $-3$ 位移（来源是 $\sum_i Q_{ii}=3$，这一点在第 \ref{sec:p3} 节还会用到）。

\textbf{定理 1.2（\file{structure-main}）}\ok\ 设 $W$ 达到 $\M$，$\sum_i p_i\delta_{v_i}$（$p_i>0$）
是其典范本征函数律，$L$ 是 $W$ 前三个特征值构成的对角阵，则：
(a) 支撑至多 $24$ 点；
(b) $W$ 在相应正测度纤维上几乎处处等于二值阶梯核 $\mathbf 1_{\{v_i\cdot v_j>0\}}$；
(c) 支撑点满足 $\sum_i p_iv_iv_i^T=I_3$、$v_i^TLv_i=\M$、$v_{i1}>0$、$v_i\cdot v_j\ne0$、
$B_{ij}=\mathbf 1_{\{v_i\cdot v_j>0\}}$、$\sum_jp_jB_{ij}v_j=Lv_i$；
(d) 对支撑中每个正点团 $J$，$\{v_iv_i^T:i\in J\}$ 线性无关，从而 $|J|\le6$。
这条断言对\emph{每一个}原始极大元成立，不是``存在某个小支撑的极大元''。

\subsection{结构线：从变分到有限支撑}

结构线的任务是回答``极大元的典范律长什么样''，
它并不直接给出常数，但它给证书线提供了那项小而关键的``正质量自对''修正。链条是四步。

\emph{第一步，二阶变分（定理 3.1，\file{negative}）。}
在\emph{实际}概率空间上做二阶变分（而不是在某个抽象参数空间上），
得到极大元处的核是半负定的。这一步是所有后续刚性结论的源头。

\emph{第二步，符号类上的秩控制（引理 3.2 \file{propagate}、引理 3.3 \file{descent}、定理 3.4 \file{finite}）。}
把支撑按坐标符号分成四个符号类，在每个类上控制正部 Gram 核的秩，
从而把典范支撑压到有限；\file{descent} 给出 $|T|\le2^{d-1}$ 的维数下降，
\file{finite} 最终给出 $|S|\le1204$。注意此处的 $1204$ 只是``有限''的定量版本，
真正的 $24$ 要靠第四步的分离观测。

\emph{第三步，矩保持变分与阈值公式（引理 2.4--2.6）。}
\file{ellipse-structure} 给 $F^TLF=\M$ a.e.；\file{row-insertion} 给 $r^TL^{-1}r\le\M$
（审核时对这一条做了独立重推，通过）；\file{contact-structure} 给阈值公式
$\sum_jp_jB_{ij}v_j=Lv_i$，这一条就是 G1 所在，现已补全。

\emph{第四步，分离二元观测与极小依赖（引理 4.1、4.2，定理 4.3，推论 4.4）。}
矩保持变分意味着典范向量对极大核的\emph{每一个二值行}都满足向量关系；
而这些行恰好能分离不同的支撑点（引理 4.2，即 G2 所在）。
再用极小依赖论证：正点团上不可能残存二阶矩关系，
于是 $\{v_iv_i^T\}$ 在正点团上线性无关，$|J|\le6$，配合四个符号类得到 $m\le24$。

需要如实说明的是，变分、克隆、卦限归约这几个机制在
Kumar--Mohar--Mojallal--Pragada 与 Kumar--Liu--Monterde--Pragada--Tait 中已有密切先例，
线性形式幂次的极小依赖的一般界属于 Li--Liu。
论文没有声称给出这类依赖的新的一般界，只声称在自己的具体假设下证明了分离观测引理，
这个自我限定在审核中被逐条验证过，名副其实。

\subsection{证书线：一个联合矩恒等式}

证书线的对象是\emph{单个}各向同性律。记
$D=\frac{3+\sqrt5}{2}-A-B$，则 \eqref{eq:lawbound} 等价于 $D\ge d$，
其中 $d=\frac{3+\sqrt5}{2}-2c=0.048687477302\ldots$。
论文的做法是同时保留四类误差量：均值亏损 $e$、径向方差 $V$、回归误差、
以及三阶与四阶的加权调和矩 $J$；然后用一个\emph{带符号的四次控制多项式}
把它们耦合进同一个精确恒等式，方向部分用 Hermite 型不等式，径向部分用比较引理，
最后整理成八项非负量之和。

具体的接力是：引理 6.1（\file{scalar}）给 $\sigma$ 的两条点态不等式；
引理 6.3（\file{radialpayment}）是``径向支付''；
命题 7.2（\file{conditional-law}）在大分支上给 $D\ge d_0$；
引理 8.2（\file{contact-square}）给 $z_*^2\le\frac32\sigma$（常数 $\frac32$ 不紧，见注 \file{rem-theta}）；
命题 8.5（\file{conditional8}）在小分支上给 $D\ge d_1$；
引理 9.1（\file{spectral-reserve}）是四次型松驰中的储备；
命题 9.3（\file{final-certificate}）把它们合成 $\M\le c$。
两个分支的分界是 $e<1/20$、$V<21/104$、$C>5/3$ 这一组阈值（式 \file{smallbranch}），
它们由 $\alpha\ge26/105$、$\beta\ge6/7$、$\rho\ge1/2$ 三个粗下界推出。
结构线在这里第二次起作用：$m\le24$ 使 $\chi-\frac{C^2}{24}\ge0$，
这是八项中最后那一项非负的依据。

谐波矩与多项式控制多项式本身都是经典工具
（Glazyrin；Der\k{e}gowska--Lewandowska；STW），
论文的贡献不在``用了控制多项式''，而在于\emph{对混合的带符号能量与绝对能量做联合控制}：
径向与方向误差不被分别丢掉，而是留在同一个恒等式里相互抵消。
这一点是改进的全部来源。

\subsection{为什么这条路线能超过 STW}

三句话概括。第一，STW 的 \eqref{eq:stw} 来自 $A\le\varphi$ 与 $B\le1$ 两次独立放缩，
而 $\gamma(3)=\varphi$ 是精确的，所以分别放缩这条路已经没有余量。
第二，本文把 $A$ 与 $B$ 放进同一个恒等式，
让``均值大''与``绝对能量大''这两种情形的误差互相支付，
于是可以在不改进任何单个投影常数的前提下改进它们的和。
第三，证书是\emph{定量的}而不是定性的：论文用的是一条精确恒等式加八项显式非负量，
没有使用任何未量化的稳定性缺口，
因此每一项的松紧都可以被计算、被指认、被改进（第 \ref{sec:slack} 节就是这么做的）。

\subsection{机器验证与数值证据}\label{sec:num}

论文所有``$>0$''与恒等式断言都在精确算术下复核：有理数用精确分数，
$\sqrt3$、$\sqrt5$ 在二次域 $\mathbb Q(\sqrt3)$、$\mathbb Q(\sqrt5)$ 中运算，全程不用浮点。
在此之上还有五组独立数值实验，全部只做交叉验证，不作为证明依据：
\begin{enumerate}[leftmargin=2em,itemsep=2pt]
\item 有限图转移在 $3\le n\le6$ 的全部 $33864$ 个标号图（$2^3+2^6+2^{10}+2^{15}$）上检查，
零反例；最小余量 $+0.0224818$，在 $K_1\vee C_5$（$n=6$）处取得。\num
\item $16$ 张符号证书的残差全部 $\le1.7\times10^{-16}$。\num
\item 对 $k=6,\ldots,20$ 的离散网格，$(A+B)/2$ 的全局最大值\emph{恰好}等于 $\rho_0$，从不超过。\num
\item $K_1\vee C_5$ 的吹起族（块大小按 $a:b=0.197271565:0.160545687$ 取整，$n=126$ 到 $2016$）
给出 $\Lam_3/n$ 从 $1.25770$ 单调升到 $1.28008$，趋于 $\rho_0$。\num
\item $4000$ 组随机各向同性律，零反例。\num
\end{enumerate}

\section{实质性进展与量化结果}\label{sec:slack}

\subsection{已经拿到手的东西}

按学位论文的口径，目前完成的是两件事，未完成的是第三件：
\textbf{A. 改进的显式上界}：$\Lam_3(G)\le cn-3$，$c<\frac{3+\sqrt5}{4}$，\ok；
\textbf{B. 结构定理}：每个极大 graphon 的典范律至多 $24$ 个原子，且原核几乎处处是二值阶梯核，\ok；
\textbf{C. 精确上界}：$\M=\rho_0$，从而 $\Lam_3(G)\le\rho_0n-3+o(n)$，\todo，这是真正的难点。
A 与 B 都依赖 G1、G2 的修补，修补已完成，因此 A、B 现在是可以写进学位论文并送审的状态。

\subsection{冠军律处的松驰归属}

改进的空间不是均匀弥散的，而是集中在两三个可以指认的位置。
把证书在冠军律（数值上取到 $\rho_0$ 的那个六原子加权轮图律）上逐项求值，
八项非负量合并展示为七组，得到下表。松驰总量 $6.210\times10^{-3}=2(c-\rho_0)$。

\begin{center}
\small
\begin{tabular}{llrl}
\toprule
项 & 含义 & 冠军律上的值 & 占松驰 \\
\midrule
$r_{\rm tr}$ & 径向迹项，$\frac89P-\E[RR'z_{\rm rad}g]$ & $3.8\times10^{-3}$ & $61.46\%$ \\
$\widehat Q$ & 带符号四次型的净贡献 & $1.6\times10^{-3}$ & $26.54\%$ \\
\quad 其中 $Q_3$ & 系数 $\frac{183}{200}-\mathrm{BIG}\,q^2$ 的最坏值放缩 & --- & $15.27\%$ \\
$\ell\, r_{\rm reg}$ & 回归误差项，$\ell=\frac{937}{1400}$ & $4.6\times10^{-4}$ & $7.44\%$ \\
$\mathcal R_0$ & 储备项 & $2.3\times10^{-4}$ & $3.68\%$ \\
$r_{\rm off}/9$ & 离对角项 & $4.9\times10^{-5}$ & $0.78\%$ \\
$K^2/1125$ & 二次储备 & $5.7\times10^{-6}$ & $0.09\%$ \\
$\frac\kappa9(\chi-\frac{C^2}{24})$ & 自对修正（用 $m\le24$） & $9.9\times10^{-7}$ & $0.02\%$ \\
\bottomrule
\end{tabular}
\end{center}

\noindent 表中``值''一列保留两位有效数字，百分比按审核报告的精确值计算，故末位可能不与两位数字直接对齐。
两条结论从这张表里直接读出来：其一，$r_{\rm tr}$ 与 $\widehat Q$ 合计占 $88\%$，
微调末四项毫无意义；其二，$\widehat Q$ 里有 $15.27\%$ 来自\emph{一处}最坏值放缩，
而这一处的系数在冠军律上是 $0.040527$、被替换成的下界是 $0.0027055$，相差 $14.98$ 倍——
这是全文最容易识别、也最应该先动的一块。

\subsection{本框架的天花板}

现有证书框架（同一条恒等式、同一组非负项、同样的分支划分）能达到的最好常数是
\begin{equation}\label{eq:ceiling}
 \frac{29+10\sqrt5}{40}=\frac{3+\sqrt5}{4}-\frac1{40}=1.2840169943749\ldots ,
\end{equation}
达到它等价于证明\emph{小分支为空}，即任何达到 $\M$ 的律都不落在
$e<1/20$、$V<21/104$、$C>5/3$ 这一区域里。若成功，
对 STW 的改进量从 $0.0243437$ 变成干净的 $\frac1{40}=0.025$，
并且 $c$ 的表达式会从 \eqref{eq:c} 那个明显是``凑出来的''数变成 \eqref{eq:ceiling} 这个有闭式意义的数，
论文的可读性与可信度都会大幅提升；两个分支给出同一个常数，本身也是自洽性的强检验。

可回收的量是 $c-\frac{29+10\sqrt5}{40}=0.000656261$，
占当前缺口 $c-\rho_0=0.003105059$ 的 $21.1\%$。
换句话说：\textbf{P1 只能回收五分之一的缺口，剩下 $78.9\%$ 在框架之外}。
这个比例必须如实告诉导师，因为它直接决定了 P1 与 P2 的性质差别。

\subsection{下界与常数项 $-3$ 的精确性}\label{sec:p3}

论文引理 2.3 用加权轮图（$K_1\vee C_5$ 的六块阶梯核，对角块取 $1$）给出 graphon 层面的
$\Sigma_3(W)\ge\rho_0=\frac{9+\sqrt5+2\sqrt{35}}{18}=1.2815681968\ldots$，
于是
\[
 \rho_0\le\M\le c .
\]
下界端点的全局最优性尚未证明。对\emph{有限图}，还需要把这个下界严格化：
设块大小 $n_0=\lfloor an\rfloor$、$n_1=\lfloor bn\rfloor$（五个），
则商矩阵 $M_n=\diag(\frac{n_i}n)(A_6+I)$ 依 $O(1/n)$ 收敛到 $M=\diag(a,b,\ldots,b)(A_6+I)$，
由 Weyl 不等式 $|\lambda_i(M_n)-\lambda_i(M)|\le\|M_n-M\|_2=O(1/n)$，
且块内均值零空间对 $\Sigma_3$ 无贡献（这正是 S1 补的那一句），
故存在图列 $G_n$ 使 $\Lam_3(G_n)/n\to\rho_0$，误差 $O(1/n)$。这一步是初等的，目前只有数值验证，\todo。

这条副线值得单独写，因为它同时给出两个有用的定位结论：
(i) 常数项 $-3$ 在这条转移路线上是\emph{精确}的，不能靠改进常数项来规避系数问题；
(ii) 上下界的差距是真实的、随 $n$ 线性增长的，不是有限 $n$ 效应
（$n=2016$ 时余量已达 $+6.2634=(c-\rho_0)n$）。

\section{当前的阻碍与风险}

以下五条是我认为必须向导师明说的，按严重程度排列。

\textbf{阻碍一：$r_{\rm tr}$ 这一项在冠军律上不消失，说明径向支付路线``付错了方向''。}
$r_{\rm tr}=\frac89P-\E[RR'z_{\rm rad}g]$，其中 $-\E[RR'z_{\rm rad}g]$ 占 $r_{\rm tr}$ 的 $87\%$，
来源是 $g(w)=\frac{1-9w^2}{25}$ 在 $|w|>\frac13$ 上为负。冠军律上 $9J-V=0.055011>0$，
所以这一项取的是正值 $3.8\times10^{-3}$ 而不是零。
一个\emph{紧}证书要求所有项在等号情形消失，因此 $r_{\rm tr}$ 根本不该出现在紧证书里。
这个判断本身就是 P2 的起点：引理 6.3 的径向支付路径很可能需要整体放弃、另找一个在 $9J=V$ 时消失的量。
对 P1 而言，它意味着``给 $r_{\rm tr}$ 一个正下界''这条子路线是真正的不确定项——
可行的小分支约束是 $e<1/20$，即 $|\E X|^2>0.95$，律必须高度集中在一个方向上，
而高度集中与 $z_{\rm rad}$ 小（$R$ 接近常值）直觉上不能同时发生；
这个``不能同时''若能量化，就是 P1 的关键引理，但目前还没有证明。

\textbf{阻碍二：$Q_3$ 的系数放缩不能简单反向操作。}
最便宜的一处改进是把 $q^2$ 保留为变量而不是先取最坏值，
但 $Q$ 对 $q^2$ 的\emph{净}依赖是符号不定的，
所以``保留 $q^2$''不等于``代入冠军值就变好''，
必须把 $Q$ 看成 $(\sqrt e,\sqrt V,\sqrt J,q)$ 的、系数含 $q^2$ 的二次型，
在小分支约束集 $\{e<1/20,\ V<21/104,\ q^2\le20/57,\ J\le V/3\}$ 上做参数化配方，
或者干脆用 SOS 重新拟合。技术上一维（$q^2$）分段有理配方可以用精确算术完成，
但工作量不小，且成功与否取决于配方后残系数能否全部保持为正。

\textbf{阻碍三：收紧小分支阈值几乎不带来收益。}
阈值 $e<1/20$、$V<21/104$、$C>5/3$ 由 $\alpha\ge26/105$、$\beta\ge6/7$、$\rho\ge1/2$ 三个粗下界推出，
用它们在 $\mathbb Q(\sqrt5)$ 中的精确值重推是零风险的机械工作，
但实测只能把 $V$ 的界收紧约 $0.16\%$。它的价值不在自身，
而在于给阻碍一、二的操作腾出空间；不能指望它单独解决问题。

\textbf{阻碍四：P2 需要的不是``调紧常数''，而是一个新的等号情形。}
八项在冠军律上的值分别是 $3.8\times10^{-3}$、$1.6\times10^{-3}$、$4.6\times10^{-4}$、
$2.3\times10^{-4}$、$4.9\times10^{-5}$、$5.7\times10^{-6}$、$9.9\times10^{-7}$，
\emph{没有一项为零}。所以 $\M=\rho_0$ 需要一个以加权轮图为等号情形的新构造。
两条候选路线：对偶/SOS 路线（冠军律的八条结构方程已全部精确知道，
包括符号模式与阈值公式，这为``猜对偶解''提供了强约束）；
结构刚性路线（用定理 1.2 归约为有限模板问题，
第一步是把 $24$ 降到 $6$——冠军律的正团结构恰是 $K_1\vee C_5$，
而 \S4 已经给出 $|J|\le6$，所以``整个支撑就是一个正团''是自然猜想，
但它与``四个符号类各 $\le6$''之间还有距离，需要新论证）。
在当前信息下 P2 不宜承诺。

\textbf{风险 R1：Ebrahimi 2008 硕士论文全文未获取。}
该论文是 $\Lam_3$ 这类问题的最早出处之一，其中可能已有相关界。
这不影响本文证明的正确性（我们改进的是 STW 的显式常数，且证明自洽），
但影响``问题归属与优先权''的表述。建议答辩前通过馆际互借获取全文核一遍。

\textbf{方法上的一条自我约束。}
极图无法由算法推出，逐阶跑算法对本问题没有意义。
本项目坚持：所有数值与优化只用于\emph{发现}常数和交叉验证，
最终证书必须用精确有理算术与二次域算术逐条验证。
这是论文中最可靠的部分，也是审核中唯一没有任何怀疑的部分。

\section{论文排版与写作建议}

以下按``先做''到``可缓''排列。前四条是我认为投稿前必须处理的，
中间部分是可读性，最后是学位论文版的版式与章节映射。
每条都给出具体位置，便于直接动手。

\subsection{T1 结构层（投稿前必做）}

\begin{enumerate}[leftmargin=2em,itemsep=4pt]
\item \textbf{在 \S1 末把证明路线图升级为编号流程。}
现有那段``Sections 2--4 prove the structural result\ldots''是散文式的，
读者要自己在脑子里画依赖图。建议改成五行带编号的 displayed 列表，
并明确标出唯一的接力点：结构线的 $m\le24$ 是证书线中 $\chi-\frac{C^2}{24}\ge0$ 的输入。
这一条最花小钱办大事，因为它直接决定审稿人对``两条线为什么必须同时存在''的理解。

\item \textbf{把命题 9.3（\file{final-certificate}）的八项做成带编号的表或 enumerate。}
定理 1.2 已按 S6 改成 (a)--(d)，同样的处理应该用在八项非负分解上：
给每项一个名字与编号（如 (T1) $r_{\rm tr}$、(T2) $\widehat Q$、\ldots、(T8) $\frac\kappa9(\chi-\frac{C^2}{24})$），
正文后续讨论时引用编号而不是复述公式。本文件 \S3.2 的那张表可以直接搬进论文，
它同时也是``改进靶子在哪''的最有力说明。

\item \textbf{附录 A 补一张常数一览表。}
现在的附录 A 用两段散文区分``唯一确定''与``选定后验证''，内容是对的，
但查阅体验差。建议改成四列表格：符号 / 数值 / 首次出现位置 / 来源类别，
把 $c$、$\cstar$、$d$、$d_0$、$d_1$、$\ell=\frac{937}{1400}$、$\eta=\frac1{10}$、$\theta=\frac6{25}$、
$\kappa=\frac3{125000}$、$\rho_0$、$41503=7^3\cdot11^2$、$z_*$、$\frac32$ 全部列入。

\item \textbf{\S10（Further questions）升格并分小标题。}
现在这一节把下界、精确值、极值图三件事混在两段里。建议拆成三个小标题：
``下界的全局最优性''、``精确值 $\M=\rho_0$''、``极值图与模板分类''，
每个小标题下先写已知、再写缺什么。学位论文的``展望''章可以直接从这里长出来。
\end{enumerate}

\subsection{T2 记号与可读性}

\begin{enumerate}[leftmargin=2em,itemsep=4pt]
\item \textbf{加一段 Notation。}
论文用到的符号至少有
\[\begin{gathered}
 \Lam_3,\ \Sigma_k,\ \M,\ A,\ B,\ D,\ d,\ d_0,\ d_1,\ c,\ \cstar,\ e,\ V,\ J,\ q,\ \sigma,\ z_*,\ \chi,\ C,\ \kappa,\ \ell,\ \eta,\ \theta,\\
 \mathcal R_0,\ r_{\rm tr},\ r_{\rm reg},\ r_{\rm off},\ K,\ P,\ g,\ B_{ij},\ L ,
\end{gathered}\]
全部散落在正文各处首次定义。建议在 \S1 末或 \S2 开头加半页 Notation 段，按``图论量 / 律量 / 证书量''三组排列。

\item \textbf{$\cstar$ 改名。}
$c_{\mathrm w}$（引理 2.3 中 $\mathcal M$ 的下界常数）与主常数 $c$ 字形太接近，
在同一段里出现时极易误读。建议改成 $w_0$ 或 $\mathcal W$。这是本文最值得做的一处记号手术。

\item \textbf{\S5.1 的长段拆三段。}
推导 $L_q-L=q_0(\zeta-x)^2(\zeta+2x)/x^2$ 的那一段现在是一整块，
建议按``恒等式 / 点态不等式 / 积分与配方''切成三段，每段一个落点。S3 补的半句正好落在第一段末。

\item \textbf{每个矩预算式加一句用途。}
式 (5.1)--(6.4)（\file{budget4}、\file{budget7}）建议逐条加一行斜体说明，
例如``此条将在命题 7.2 中用于吸收 $\widehat Q$ 的三次部分''。
审核时我自己就多次需要回翻才想起某条预算的用途，审稿人更会如此。

\item \textbf{主要引理后加一行``用途''注。}
共约八处：引理 2.2（$\Sigma_2(W)\le7/6$）、2.5、2.6、3.3、4.2、6.1、6.3、8.2。
每处一句话，说明它在哪一步被调用、调用的目的是什么。
这类``路标注释''对长证明的可读性提升远超任何版式调整。
\end{enumerate}

\subsection{T3 严谨性表述（小，但审稿人会抓）}

\begin{enumerate}[leftmargin=2em,itemsep=4pt]
\item \textbf{附录 A 的 $33864$ 要写明口径。}
现在写的是``all $33864$ graphs of order $3\le n\le6$''。
$33864=2^3+2^6+2^{10}+2^{15}$ 是\emph{标号}图（即全部邻接矩阵）的数目；
未标号图只有 $4+11+34+156=205$ 个。建议改成
``all $33864$ labelled graphs ($2^3+2^6+2^{10}+2^{15}$) of order $3\le n\le6$''。
不改的话，任何一个动手核对的审稿人都会在这里卡住并怀疑其余数字。

\item \textbf{摘要里补上 $|J|\le6$。}
摘要现在只说``at most 24 support points''，但真正被证书使用的是正团上的 $|J|\le6$
（它给出 $\chi-\frac{C^2}{24}\ge0$）。建议摘要同时提这两条。

\item \textbf{定理 1.1 的假设前移。}
``Only a finite second moment is assumed; zero vectors and unbounded radii are allowed''
现在写在结论之后，建议移到陈述开头作为假设行，
这样``对所有二阶矩有限的律''这个范围限定第一眼就能看到。

\item \textbf{非标准记号给前向引用。}
$\mathrm{BIG}$、$\widehat Q$、$z_*$、$\chi$ 这类记号在首次出现处应注明定义所在的式号或引理号，
尤其 $\mathrm{BIG}$ 是一个具体有理数，读者不该需要往后翻两页才知道它是什么。
\end{enumerate}

\subsection{T4 版式（LaTeX 层面）}

\begin{enumerate}[leftmargin=2em,itemsep=4pt]
\item 现在是 \file{article}、11pt、\file{margin=28mm}，投稿版合适，保留。
学位论文版按学校模板改（一般 12pt、上下 2.5cm、左侧 3cm 装订线），并加 \file{setspace} 的 \file{onehalfspacing}。
\item \file{hyperref} 现为 \file{hidelinks}；学位论文版建议开启彩色链接，便于导师在 PDF 里跳转（本文件即如此设置）。
\item 补 MSC 2020 分类号（建议 05C50、15A18、52C45）；keywords 已有，保留。
\item \textbf{参考文献格式统一。}现在 DOI、arXiv 号、期刊卷页三种混用；
2026 年的预印本版本号也不统一（STW 标 v1，KLMPT 标 v2，其余 v1）。
建议统一为``作者，标题，期刊 \emph{卷} (年) 页，DOI''或``arXiv:编号 (年)''两种模板之一，
并对预印本统一标注所引版本。
\item 建议引入 \file{cleveref}，把正文中手写的 \file{Theorem~\textbackslash ref\{\}} 与 \file{\textbackslash eqref\{\}} 统一交给宏包，
减少交叉引用风格漂移。
\item \textbf{加两张图。}论文目前完全没有插图，这在学位论文里是明显的短板。建议：
(i) 冠军律六个支撑点在单位球面上的位置，标出正团结构 $K_1\vee C_5$ 与符号类划分；
(ii) 吹起族 $\Lam_3/n$ 随 $n$ 的收敛曲线（$n=126$ 到 $2016$），
同图画出 $c$ 与框架天花板 $\frac{29+10\sqrt5}{40}$ 两条水平线。
第二张图在答辩现场比任何公式都更能说明``缺口是真实的、随 $n$ 线性增长''。
\end{enumerate}

\subsection{T5 学位论文化的章节映射}

期刊式短文与学位论文的信息密度要求不同。建议映射如下：
\begin{center}
\small
\begin{tabular}{ll}
\toprule
学位论文章节 & 来源 \\
\midrule
第 3 章\ 预备知识（graphon、投影常数、STW 归约） & \S1--\S2 前半 \\
第 4 章\ 极大元的结构（结构线） & \S2--\S4 \\
第 5 章\ 联合矩证书（证书线） & \S5--\S9 \\
第 6 章\ 主定理与有限图转移 & \S9 末、\S10 前 \\
第 7 章\ 机器验证与数值实验 & 附录 A $+$ 本文件 \S2.6 \\
第 8 章\ 展望：精确上界与下界 & \S10 $+$ 本文件 \S3.3--\S3.4 \\
\bottomrule
\end{tabular}
\end{center}
另外，学位论文需要一段``本文贡献与已有工作的区别''，
可直接把论文 \S1.1（Method and relation to earlier work）扩写。
其中的自我限定——不声称给出线性形式幂次极小依赖的新一般界、
不声称 $c$ 最优、不声称改进了所有涉及额外图参数的估计——
建议原样保留，答辩时这些限定是加分项而不是减分项。

\section{逐页讲稿（自编译 26 页版）}

幻灯片为自编译版本：不依赖任何模板 skill，版式由 \file{python-pptx} 直接排，
公式与三张图表由 \file{matplotlib} 渲染成 PNG 嵌入，PDF 经 WPS COM 导出。
共 26 页，全套讲完约 20--25 分钟，留出提问时间。
下面每页给``要记住什么''与口播稿，重点页附预设提问与答法。

\subsection*{第 1 页　封面}
要记住：题目、四个关键数字、四个状态标签。口播：这次汇报的是毕业论文主线的一个阶段结果；
四个数字卡是全套的索引——现在在哪、改进了多少、这条路的终点、真值大概在哪。

\subsection*{第 2 页　一页结论}
要记住：``已经拿到手的''与``还没有拿到手的''两栏，以及右栏的缺口诊断。
口播：左上是 A（改进上界）与 B（结构定理）两件已完成的事和审核状态；左下是三件未完成的事；
右栏预告缺口集中在 $r_{\rm tr}$ 与 $\widehat Q$ 两处。
预设提问：\emph{天花板与上确界有什么区别？}答：天花板是``用现在这条证书还能达到的最好常数''，
上确界是数值上观察到的真实极值、由显式加权轮图构造给出，是下界。

\subsection*{第 3 页　汇报路径}
六个板块。15 秒带过，不逐条念。

\subsection*{第 4 页　背景与桥}
要记住：本页给只读过 Nikiforov 图极限的听众搭桥。口播：先回顾 graphon 底座（阶梯 graphon、
blow-up 不变、割度量紧性），指出 $\lambda_2,\lambda_3$ 不是同态密度故不能直接套；再给出本文
第一块砖——Ky Fan 把 $\Lambda_3$ 写成迹型泛函，取极限后结构被一个各向同性律 $\mu$ 捕获，
谱问题变成矩问题。预设提问：为什么迹型就能转移？答：Ky Fan 是有限维 max-trace，graphon 极限下
迹型泛函上半连续且可由三维子空间见证，故可转移（毕业论文 Lemma 2.3）。

\subsection*{第 5 页　问题设定}
要记住：$\Lam_3$ 保留符号、只以 $n$ 为约束；graphon 语言与``补零''约定只用于 $\Sigma_k$。
口播：出处是 Ebrahimi 2008 与 Mohar 2009；右栏给出归一化语言 $A$、$B$ 与它们的单独上界，
并点明 $k=3$ 为什么最有信息量：$\gamma(3)=\varphi$ 精确。

\subsection*{第 6 页　已知基准 STW}
要记住：STW 的界是两次独立放缩之和，而 $\varphi$ 精确，所以分别放缩没有余量。
口播：混合投影常数问题 STW 自己在其 \S6.2 提出但没解决，那就是本文的切入点；
右栏一句话：改进必须是结构性的。

\subsection*{第 7 页　主定理}
要记住：$c$ 是显式代数数，不是拟合结果；两条等价表述；$-3$ 的来源。
口播：律版是主体，图版是转移；转移产生精确的 $-3$（来源 $\sum_iQ_{ii}=3$）。
预设提问：\emph{$-3$ 能不能改进？}答：在这条转移路线上它是精确的，所以系数问题不能靠常数项回避。

\subsection*{第 8 页　证明架构}
要记住：两条线、一个接力点。口播：结构线负责收敛、证书线负责闭合；
唯一接力点是 $m\le24$ 撑起 $\chi-\frac{C^2}{24}\ge0$。底部一行交代先例与自我限定。
预设提问：\emph{结构线能不能不用？}答：可以，但第八项要换更粗的界，常数会变差；
只有结构线则完全给不出常数，二者缺一不可。

\subsection*{第 9 页　结构线 I}
要记住：三步链条与``1204 只是定量的有限''。口播：二阶变分在\emph{实际}概率空间上做；
表格里三个行级引理，其中 2.6 就是 G1 所在，已补约十行。

\subsection*{第 10 页　结构线 II}
要记住：定理 1.2 的 (a)--(d)，以及从 1204 到 24 的三步。
口播：右栏强调承重性——$m\le24$ 撑起证书第七项；G2 就是引理 4.2 补的那一句。

\subsection*{第 11 页　证书线 I}
要记住：保留四类误差、一个联合恒等式、七组非负项的编号表。
口播：控制多项式 $f$ 的五个系数由插值唯一确定，所以 $41503=7^3\cdot11^2$ 不是凑的；
表 T1--T7 是后面所有讨论的公共语言。

\subsection*{第 12 页　证书线 II}
要记住：七步接力链与分支阈值；两处非紧常数如实标注。
口播：左栏是``已知不紧但不影响''的两处（注 \file{rem-theta}、\file{rem-radii}），
右栏是设计意图——误差不被分别丢掉，而是在同一恒等式里互相抵消。

\subsection*{第 13 页　为什么超过 STW}
三条理由逐条念，第三条引出``缺口可诊断''，是下一页的入口。

\subsection*{第 14 页　机器验证}
要记住：六组实验、完备度口径、以及 $33864$ 是标号图。
口播：数值只交叉验证不作证明；完备度 100 与 90 的区别要讲清楚。
预设提问：\emph{$33864$ 怎么来的？}答：$2^3+2^6+2^{10}+2^{15}$；未标号只有 205 个，
论文附录建议写明 labelled，已列入排版建议 T3。

\subsection*{第 15 页　松驰归属}
要记住：$88\%$ 集中在两项，其中 $15.27\%$ 是一处单一放缩。口播配合条形图讲右栏两条结论。

\subsection*{第 16 页　$Q_3$ 的 15 倍}
要记住：四个数字卡与``为什么不能直接代冠军值''。
口播：$Q$ 对 $q^2$ 的净依赖符号不定，所以正确做法是参数化配方或 SOS 重拟合；这是 P1 第一条。
预设提问：\emph{把这处放缩去掉不就行了？}答：去掉最坏值不等于变好，必须整体重新配方并验证残系数全为正。

\subsection*{第 17 页　常数阶梯}
要记住：三段占比 88.7 / 2.4 / 8.9。口播：瀑布图的形状本身就是结论——
最难的部分已经过去，但最后一公里需要换工具。
预设提问：\emph{为什么不直接做 P2？}答：P2 需要逐项取零的新证书，目前八项无一为零；
P1 属重新配方，成功率高且能把 $c$ 变成干净闭式。

\subsection*{第 18 页　框架天花板}
要记住：$\frac{29+10\sqrt5}{40}$ 与达到它的充要条件（小分支为空）。
口播：成功后改进量恰为 $\frac1{40}$；但必须如实说只能回收 $21.1\%$，其余 $78.9\%$ 在框架外。
预设提问：\emph{P1 把握多大？}答：三条子路线里两条信心较高、一条真正不确定，综合``有实质机会但不稳''。

\subsection*{第 19 页　审核 G1/G2}
要记住：两处缺口的``缺什么''与``怎么补''，以及各自的承重位置。
口播：审核是独立做的，33 页；补完后无连带影响。
预设提问：\emph{为什么这两处算承重？}答：G1 撑阈值公式即 $A+B=2\M$ 的入口，G2 撑 $m\le24$ 即证书第七项。

\subsection*{第 20 页　S1--S7 与附录 A}
表格逐行 10 秒带过；强调效果：19 页变 20 页、编译无警告，附录 A 直接回应``系数是调出来的''质疑。

\subsection*{第 21 页　阻碍清单}
五条逐条念，不美化。底部一行是方法上的自我约束：算法只用于发现与交叉验证。

\subsection*{第 22 页　优先级矩阵}
要记住：右下两条是\emph{明确排除}的，不是没想过。
预设提问：\emph{P3 收益低为什么还要做？}答：它给出``$-3$ 精确''与``缺口随 $n$ 线性增长''两个定位结论，
论证 P1/P2 的必要性，且几乎零风险。

\subsection*{第 23 页　P0--P3 路线图}
要记住：排序依据四条与决策点。口播：把握度是主观值不是完成度；
先花两天做零风险的阈值重推并画可行域，再决定 P1 第 2 条投入多少。
预设提问：\emph{P1 失败怎么办？}答：论文仍完整——A 与 B 已站住；失败过程给出的可行域刻画本身可写成一节。

\subsection*{第 24 页　结论与排版建议}
左栏三句话收尾数学部分；右栏 T1--T5 表告诉导师``接下来论文怎么改''，
并说明 T1/T2 决定可读性、T3 是口径问题、T4/T5 决定能否长成一章学位论文。

\subsection*{第 25 页　适用范围与披露}
要记住：口径与论文一致，主动披露。口播：加权 $[0,1]$ 是适用范围（Prop 6.8/6.10），不是新定理；
加权形式不新颖（STW Prop 6.1 证明本身只用 $a_{ii}=0$、$0\le a_{ij}\le1$；Mohar 2009 已覆盖对称矩阵）；
下界常数 $c_w$ 与 $K_1\vee C_5$ 模板已被 Zenodo 自存档条件性先提出，论文注 3.2 已让渡；
主定理与 24-原子结构二次查新未被抢先。预设提问：既然加权不新颖为何还写？答：作为适用范围与
否定性结果（无渐近增益）留存，并在 Declaration 如实披露。

\subsection*{第 26 页　结语}
口播：主攻小分支为空，副线吹起族严格渐近；方法上数学推理优先、算法只作交叉验证。
右栏请导师重点指正三处：P1 投入比例、P2 定位、学位论文章节映射。

\section{后续优先级与时间预算}

排序原则保持不变：\textbf{数学方法优先，算法只用于发现与交叉验证；
每条路线都要有明确的收益量化与可验证的判据；先做零风险的，再决定是否 all-in。}

\subsection*{P0\ 论文加固\quad（已完成）}
G1、G2 两处修补 $+$ S1--S7 $+$ 附录 A，论文重新编译通过（19 页 $\to$ 20 页）。
本轮另加：排版建议清单（本文件 \S5）、汇报幻灯片与讲稿。
剩余可选动作：把 \S5 的 T1 四条与 T3 第 10 条（$33864$ 口径）在投稿前落实，预计半天。

\subsection*{P1\ 证明小分支为空\quad（主攻，中等难度）}
目标：$c=\frac{29+10\sqrt5}{40}$，改进量恰为 $\frac1{40}=0.025$，回收当前缺口的 $21.1\%$。
三条子路线按优先级：
\begin{enumerate}[leftmargin=2em,itemsep=3pt]
\item \textbf{保留 $q^2$ 做参数化配方}（信心较高）。把 $Q$ 看成
$(\sqrt e,\sqrt V,\sqrt J,q)$ 的、系数含 $q^2$ 的二次型，
在小分支约束集上做一维分段有理配方，用精确算术验证残系数全为正。
针对的是 $15.27\%$ 的松驰与 $14.98$ 倍的系数差距。
\item \textbf{给 $r_{\rm tr}$ 一个正下界}（真正的不确定项）。需要证明
``小分支迫使律高度集中（$e<1/20$ 即 $|\E X|^2>0.95$）''与
``$z_{\rm rad}$ 小（$R$ 接近常值）''不能同时发生。
建议先专门研究小分支下 $(e,V,J,q^2)$ 的联合可行域并画图，
用图形判断这条是否有希望，再决定投入。
\item \textbf{收紧小分支阈值}（零风险，但收益极小）。用 $\alpha,\beta,\rho$ 在 $\mathbb Q(\sqrt5)$ 中的精确值
重推 $e<1/20$、$V<21/104$、$C>5/3$；实测 $V$ 只收紧约 $0.16\%$，
价值在于给第 1、2 条腾操作空间。
\end{enumerate}
建议节奏：先花两天做第 3 条并画出可行域，据此决定第 2 条投入多少；第 1 条与第 3 条并行。

\subsection*{P2\ 精确上界 $\M=\rho_0$\quad（长期目标，高难度，不宜承诺）}
需要一个在冠军律上\emph{逐项取零}的证书；目前八项无一为零。
两条候选路线：
(i) 对偶/SOS——冠军律的八条结构方程已全部精确知道（含符号模式与阈值公式），
为构造对偶可行解与互补松弛提供了强约束；
由 $9J-V=0.055\ne0$ 可知 $r_{\rm tr}$ 不该出现在紧证书里，
即引理 6.3 的径向支付路径很可能要整体放弃；
(ii) 结构刚性——用定理 1.2 归约为有限模板问题，
\emph{第一步是把 $24$ 降到 $6$}（冠军律的正团结构恰为 $K_1\vee C_5$，
\S4 已给 $|J|\le6$，故``整个支撑就是一个正团''是自然猜想，但尚需新论证）。
这一步对两条路线都有用，建议作为 P2 的唯一近期动作。

\subsection*{P3\ 下界的严格渐近\quad（副线，低风险，建议完成）}
把 $\rho_0$ 从 graphon 层面升级为有限图列：块大小取整 $+$ Weyl 不等式给 $O(1/n)$ 误差，
配合 S1 补的``块均值零子空间 $\Rightarrow\Sigma_3$ 等于商谱''即可写成完整一节。
产出两个定位结论：常数项 $-3$ 精确、上下界差距随 $n$ 线性增长。建议与 P1 并行。

\section{交付物与复现}

\subsection*{文件清单}
\begin{center}
\small
\begin{tabular}{@{}p{8.1cm}p{5.6cm}@{}}
\toprule
文件 & 说明 \\
\midrule
\file{D:\textbackslash top3\_joint\_moment\_bound.tex/.pdf} & 论文正文，22 页纠正版（期刊）；毕业论文 thesis\_main/main.tex，60 页 Master's \\
\file{审核报告\_top3\_joint\_moment\_bound.pdf} & 33 页独立审核报告（上一轮产出） \\
\file{汇报\_前三大特征值上界\_20261102/讲稿与排版建议.pdf} & 本文件 \\
\file{\ldots/汇报PPT\_自编译\_前三大特征值上界\_20261102.pptx} & 26 页自编译幻灯片（可编辑） \\
\file{\ldots/汇报PPT\_自编译\_前三大特征值上界\_20261102.pdf} & 同上的 PDF 版（WPS COM 导出） \\
\file{\ldots/assets/} & 公式与三张图表的 PNG 素材 \\
\file{\ldots/\_dashi\_弃用/} & 上一版 dashi 模板产物，已弃用，仅留档 \\
\bottomrule
\end{tabular}
\end{center}

\subsection*{复现命令}
论文与本文件用本地 MiKTeX 编译：
\begin{quote}\small\raggedright
\file{pdflatex top3\_joint\_moment\_bound.tex}（跑两遍，解交叉引用）\\
\file{xelatex 讲稿与排版建议.tex}（跑两遍，中文用 \file{ctexart}）
\end{quote}
幻灯片（自编译链，不依赖任何模板 skill）：
\begin{quote}\small\raggedright
\file{C:\textbackslash Python314\textbackslash python.exe -X utf8 -u \_mk\_assets.py}（matplotlib 渲染公式与三张图表到 \file{assets/}）\\
\file{<bundled python> \_mk\_pptx.py}（\file{python-pptx} 版式引擎，26 页；\file{\_ppt\_lib.py} 为版式原语，\file{\_ppt\_slides.py} 为页面内容）\\
PDF：WPS COM \file{KWPP.Application} 的 \file{SaveAs(..., 32)}
\end{quote}
说明：上一版 dashi 模板路线在本机因缺少 \file{openssl} 无法走官方 PDF 导出（曾以仓库根的
\file{\_export\_pdf.mjs} 绕过）；该路线已整体弃用，自编译链的 PDF 导出改走 WPS COM，无额外依赖。
全部数值实验脚本以 \file{\_aud*.py}、\file{\_p0\_patch.py} 等前缀保存在仓库根目录，
运行方式为 \file{C:\textbackslash Python314\textbackslash python.exe -X utf8 -u <脚本>}。

\vspace{1em}
\noindent\textbf{说明。}本文件中所有数字均来自论文正文、附录 A 与 33 页审核报告的实测记录，
标注 \ok\ 者为已证明的结论，标注 \num\ 者为数值实验结果（不作为证明依据），
标注 \todo\ 者为尚未完成、但路线已明确的部分。
\end{document}
